\documentclass{article}
\usepackage{algorithm, algorithmic, amsfonts, amsmath}

\begin{document}
\begin{algorithm}[!h]
\renewcommand{\algorithmicrequire}{\textbf{Input:}} % show \REQUIRE as Input
\renewcommand{\algorithmicensure}{\textbf{Output:}} % show \ENSURE as Output
	\caption{Partial-OT PU (PU learning by partial optimal transport)} % algorithm title
	\begin{algorithmic}[1] % [1] numbers every line
	    \REQUIRE
	    Unlabeled (source) samples $(\mathcal{X},p)$; positive (target) samples $(\mathcal{Y},q)$;
	    class prior $\pi$; initial plan $T^{(0)}$
	    \ENSURE
	    Transport plan $T$; the transported unlabeled points are identified as positive
	    \STATE \textbf{Formulation:} transport a mass $s=\pi$ from the unlabeled set to the positive set % only part of the mass is transported, which is the PU structure: a fraction pi of U is positive
	    \STATE Compute cost matrices $C^{s}$ and $C^{t}$;
	    build $\bar{p}=[\,p,\ \|q\|_{1}-s\,]$ and $\bar{q}=[\,q,\ \|p\|_{1}-s\,]$ % a dummy point absorbs the untransported mass, turning the partial problem into a balanced one
	    \FOR{$k=0,1,2,3,\dots$}
	    \STATE $G^{(k)}\leftarrow M(C^{s},C^{t})\circ T^{(k)}$
	    \hfill \emph{gradient} $\nabla J_{C^{s},C^{t}}(T^{(k)})$
	    \STATE $\bar{T}^{(k)}\leftarrow \arg\min_{T\in\Pi(\bar{p},\bar{q})}\langle \bar{G}^{(k)},T\rangle_{F}$
	    \hfill \emph{solve a partial-W problem} % the linear subproblem of Frank-Wolfe is itself a partial Wasserstein problem
	    \STATE Get $\tilde{T}^{(k)}$ from $\bar{T}^{(k)}$
	    \hfill \emph{remove last row and column} % drop the dummy point
	    \STATE Compute $\gamma^{(k)}$ by line search
	    \STATE $T^{(k+1)}\leftarrow (1-\gamma^{(k)})T^{(k)}+\gamma^{(k)}\tilde{T}^{(k)}$
	    \hfill \emph{update}
	    \ENDFOR
	    \RETURN $T^{(k)}$;\quad label $x_{i}\in\mathcal{X}$ positive iff its row of $T$ carries mass % an unlabeled instance whose mass is transported is declared positive
	\end{algorithmic}
\end{algorithm}
% Chapel, Alaya, Gasso, NeurIPS 2020, "Partial Optimal Transport with Applications on
% Positive-Unlabeled Learning". Transcribed line by line from Algorithm 1 of the original, the Frank-Wolfe solver for
% partial-GW, with the PU modelling of line 1 and the decision rule of the last line added from Section 4.2.
% Modelling: classical optimal transport requires the two distributions to carry equal mass and all of it to move.
% The PU structure breaks both, since only a fraction pi of the unlabeled set is positive. Partial optimal transport
% names the transported mass s = pi explicitly, so which unlabeled instances are moved onto the positive distribution
% is itself the identification of the positives, with no classifier to train afterwards.
% Difference from others in the branch: Self-PU, PULNS and Robust-PU all select gradually during training, whereas
% here selection is a transport problem solved directly, and the solution is the selection.
% Implementation note: the original states that the algorithm has been merged into the Python Optimal Transport (POT) library.
\end{document}
